\documentclass[12pt]{article}
\usepackage[a4paper, margin=1in]{geometry}
\usepackage{amsmath, minted, enumitem, cmbright}
\renewcommand{\thesection}{\Alph{section}}
\renewcommand{\thesubsection}{\thesection.\arabic{subsection}}
\renewcommand{\t}{\texttt}
\begin{document}
\section{Sorting Problem}
\subsection{Manual Test Cases}
\begin{enumerate}[label=\roman*.]
\item Answer: \t{1 2 3 4 5 6 7}

Explanation: Already in sorted order.
\item Answer: \t{7 6 5 4 3 2 1}

Explanation: Already in reverse sorted order.
\item Answer: \t{2 7 5 3 1 4 6}

Explanation: \t{20000} is the 2nd smallest, \t{70000} is the 7th smallest, etc.
\item Answer: \t{2 7 5 4 6 3 1}

Explanation: \t{-5} is the 2nd smallest, \t{12} is the 7th smallest, etc.
\end{enumerate}
\subsection{Propose an Algorithm and Analyze It}
As the problem stated, there are at most $2\times 10^5$ distinct integers in this range, same as the largest $n$. Thus we can solve this problem via counting sort in $O(n)$ time.

The idea is as follows:
\begin{itemize}
\item Create $2\times 10^5$ boolean `buckets' as \t{False} initially.
\item Traverse the given list; for each integer, set the corresponding bucket to \t{True}.
\item Traverse the buckets from smallest to largest; for each \t{True} value, assign the corresponding integer the next rank ($1,2,3,\dots$).
\item Traverse the original list; for each integer, print its assigned rank.
\end{itemize}

A Python code is given below.
\begin{minted}{python}
class Solution:
    def sortingProblem(self, n: int, A: List[int]) -> List[int]:
        a = [False] * 200000
        r = [0] * 200000
        for i in A:
            a[i + 100000] = True
        rank = 1
        for i in range(-100000, 100000):
            if a[i + 100000]:
                r[i + 100000] = rank
                rank += 1
        return " ".join([str(r[i + 100000]) for i in A])


tests = [[4, [777, -77, -7, 7777]],
         [7, [1, 2, 3, 4, 5, 6, 7]],
         [7, [7, 6, 5, 4, 3, 2, 1]],
         [7, [20000, 70000, 50000, 30000, 10000, 40000, 60000]],
         [7, [-5, 12, 7, -2, 9, -3, -100000]]]
for t in tests:
    print(f"{t} -> {Solution().sortingProblem(*t)}")
'''
[4, [777, -77, -7, 7777]] -> 3 1 2 4
[7, [1, 2, 3, 4, 5, 6, 7]] -> 1 2 3 4 5 6 7
[7, [7, 6, 5, 4, 3, 2, 1]] -> 7 6 5 4 3 2 1
[7, [20000, 70000, 50000, 30000, 10000, 40000, 60000]] -> 2 7 5 3 1 4 6
[7, [-5, 12, 7, -2, 9, -3, -100000]] -> 2 7 5 4 6 3 1
'''
\end{minted}
\section{Stack with `findMin()' Operation}
\textit{Note that the original problem is written in Java but I will do this in Python.}

The idea is to maintain a second stack to keep track of the current minimum so far.

A Python code is given below.
\begin{minted}{python}
class SpecialStack:
    def __init__(self):
        self.s = []
        self.m = []

    def peek(self) -> int:
        return self.s[-1] if self.s else -1

    def pop(self) -> int:
        if not self.s:
            return -1
        if self.m:
            self.m.pop()
        return self.s.pop()

    def push(self, value: int) -> None:
        self.s.append(value)
        if self.m:
            self.m.append(min(value, self.m[-1]))
        else:
            self.m.append(value)

    def findMin(self) -> int:
        return self.m[-1] if self.m else -1


if __name__ == "__main__":
    S = SpecialStack()
    assert S.peek() == -1
    assert S.findMin() == -1
    for x in [2, 7, 8, 6, 1, 5, 1]:
        S.push(x)
    assert S.peek() == 1
    assert S.findMin() == 1
    assert S.pop() == 1
    assert S.peek() == 5
    assert S.findMin() == 1
    assert S.pop() == 5
    assert S.pop() == 1
    assert S.peek() == 6
    assert S.findMin() == 2
    print("Test cases are all executed correctly.")
\end{minted}
\section{Adding $n$ Integers}
The idea is to use a greedy algorithm: just take the two smallest numbers, pay a cost of their sum, then push their sum back, repeat until there is only one element left and return the total cost.
\subsection{Manual Test Cases}
\begin{enumerate}[label=\roman*.]
\item Answer: \t{10}

Explanation: $3+7=10$.
\item Answer: \t{24}

Explanation: Merge every pair of $1$'s together first, giving a cost of $2\times4=8$. Then merge every pair of $2$'s, giving $4\times 2=8$. Then merge the two $4$'s giving $4+4=8$. Thus the cost is $8+8+8=24$ in total.
\item Answer: \t{960}

Explanation: Do $32+32=64$, then $64+64=128$, then $128+128=256$, then $256+256=512$. Thus the cost is $64+128+256+512=960$ in total.
\item Answer: \t{35}

Explanation: Do $1+2=3$, then $3+3=6$, then $5+5=10$, then $6+10=16$. Thus the cost is $3+6+10+16=35$ in total.
\end{enumerate}
\subsection{Special Case 1}
Simply return $n\log_2{n}$. As $n$ is a power of $2$, in Python we can do
\begin{minted}{python}
return n * (n.bit_length() - 1)
\end{minted}
which has time complexity $O(\log n)$ as \t{bit\_length} is $O(\log n)$.
\subsection{Special Case 2}
We can always take the first two elements and merge them, repeat until only one element is left. A Python code is given below.
\begin{minted}{python}
s = A[0]
cost = 0
for i in A[1:]:
    s += i
    cost += s
return cost
\end{minted}
\subsection{Propose an Algorithm and Analyze It}
Maintain a min-heap to get the smallest values efficiently in $O(\log n)$. Thus the time complexity is $O(n\log n)$. A Python code is given below.
\begin{minted}{python}
class Solution:
    def addingnIntegers(self, n: int, A: List[int]) -> int:
        from heapq import heapify, heappush, heappop
        heapify(A)
        cost = 0
        while len(A) > 1:
            s = heappop(A) + heappop(A)
            cost += s
            heappush(A, s)
        return cost


tests = [[3, [1, 2, 4]],
         [2, [7, 3]],
         [8, [1, 1, 1, 1, 1, 1, 1, 1]],
         [5, [32, 256, 128, 64, 32]],
         [5, [5, 1, 3, 2, 5]]]
for t in tests:
    print(f"{t} -> {Solution().addingnIntegers(*t)}")
'''
[3, [1, 2, 4]] -> 10
[2, [7, 3]] -> 10
[8, [1, 1, 1, 1, 1, 1, 1, 1]] -> 24
[5, [32, 256, 128, 64, 32]] -> 960
[5, [5, 1, 3, 2, 5]] -> 35
'''
\end{minted}
\end{document}